% 【研读实验室】所属：第4章（由 chapters/revised/04-sources.tex \input 引入）
\minitopic{研读实验室：记忆失败发生在哪一层}\index{记忆保存论}\index{来源监控}\index{内省}\index{归纳} “记错了”至少有四种机制。第一，编码不足：事件发生时主体未注意；第二，保存变化：信息随时间衰退或被重构；第三，提取失败：信息仍有痕迹却暂时想不起；第四，来源监控错误：记得内容，却把听来的事当成亲历。将它们统称记忆不可靠，既不能解释个案，也不能指导核验。 记忆还可保存命题、经验形象、技能与个人事件。会骑车不依赖持续记住学习时的每个命题；知道首都名称可能在忘记教师和课堂后仍保留。讨论“记忆是否只是保存知识”时，应标明对象：命题记忆的确证可能沿时间保存，情景记忆则包含自我置身过去的体验结构。 保存论认为，记忆主要维持先前获得的知识或确证；生成论强调，当前记忆状态本身可在整合痕迹后提供新的认识地位。若主体最初草率听见消息，后来只记得内容，保存论难以让不良起点变好；若多段不完整记忆在后来形成合理整体，生成论较灵活，却要防止融贯重构制造虚假确定性。

\minitopic{遗忘来源与认识地位}我们常记得“巴黎是法国首都”，却忘记从何学来。强内在主义若要求主体当前能给出来源，许多普通长期知识会消失；外在主义可说原学习链可靠且记忆正常保存，来源遗忘不取消知识。中间立场要求主体至少没有具体击败理由，并能在重要情形下查证。 来源遗忘的风险不均。基础常识由许多独立学习和实践反复支持，即使单一来源消失也不脆弱；冷门指控可能只来自一次不可靠转述，忘记来源会让主体高估地位。合理评价应考虑内容是否有多重支持、原渠道是否大体可靠、后来是否被独立使用和校正。 “我记得$p$”也可能是一种高阶证据。记忆呈现通常给予初步相信权利，但若主体知道自己刚看过大量相似假信息、经历引导性提问或有特定记忆障碍，权利被削弱。一般“记忆有时会错”不取消每项记忆，正如视觉可错不取消所有观看。 记忆信心不是准确度的透明指标。反复讲述、他人确认和情绪强度会提高鲜明度，却可能同时引入重构。核验时要优先最早记录、独立证人、时间戳和能区分竞争叙事的细节，而非简单多数或最自信者。

\minitopic{内省与自我知识的分层}主体对当前疼痛通常具有特殊权威，但这不等于关于疼痛的每项判断都不可错。可以区分：正在有某种不适体验；把它归类为疼痛或焦虑；判断原因是胃病；预测它会持续。越向解释和预测移动，越依赖概念、背景知识与外部证据。 自我知识理论包括内在扫描、表达主义和透明性路线。扫描模型把内省类比观察内部状态，面临内部对象与识别机制问题；表达主义认为真诚说“我相信$p$”往往直接表达信念，而非报告一次内观发现；透明性路线通过考虑世界是否支持$p$来回答自己是否相信$p$。三者可能解释不同状态，不必强求单一机制。 自我欺骗说明第一人称权威可受动机影响。主体可能准确知道自己紧张，却不承认嫉妒；也可能口头接受某价值，行为却持续显示相反倾向。第三人称行为证据不能自动压过当事人，但当行为跨情境稳定、当事人解释不断临时变化时，它构成相关挑战。 心理治疗、日记和他人反馈可增加自知，说明自我知识并非完全私有。社会解释也可能压迫性地替主体定义经验，因此需要双重约束：认真对待第一人称呈现，同时允许可说明的外部证据纠正分类与因果判断。特殊权威不是绝对否决权。

\minitopic{先验确证的边界}先验不是“出生时就知道”，也不是“完全不需要任何经验”。理解语言、学习符号和形成概念通常依赖经验；先验主张是，某命题的最终确证不依赖特定经验事实。理解“所有单身者未婚”后，主体不需调查每位单身者，但仍需学会相关词语。 分析真理、逻辑推理与数学证明常被视为先验候选。它们的认识问题不同。分析真理依赖意义关系；逻辑知识涉及有效规则如何被认识；数学证明需要对长推导的记忆与检查。即使结论先验，实际获得过程仍可能因抄写错误或能力缺陷而失败。 理性直觉在哲学中常提供“似乎必然如此”的呈现。若把任何强烈直觉都当作先验证据，会放过文化偏见与理论训练效应；若完全拒绝，又难以解释我们如何评估简单逻辑和反例。较谨慎立场把直觉视为可被击败的初步证据，并通过论证、案例变体和跨主体检验校准。 先验与必然需要持续区分。一个命题如何被确证是认识分类，它在其他可能情形是否为真是模态分类。把两者混同，会错误地从“我不用观察便知道”推出“事情不可能不同”，或从“命题必然真”推出“任何相信方式都产生知识”。

\minitopic{归纳：证明、辩护与改进}休谟问题指出，从过去观察到未来或未观察实例的推移不能由演绎保证；用“归纳过去成功”辩护归纳又使用了归纳\parencite{hume2000}。回应首先要说明目标：寻找非循环演绎证明、说明归纳是理性的规范，还是设计表现更好的具体方法。目标不同，循环的严重性与所需答案不同。 实用主义回应可能说，无论世界是否规则，若有可利用规律，归纳方法有机会发现；不用归纳则没有更好策略。这给出策略辩护，不证明每次归纳结论为真。贝叶斯路线用概率一致性和更新规则刻画信念变化，却仍需先验概率和似然模型，不能单靠公式消除经验假设。 自然主义研究哪些归纳方法在何种环境校准。随机抽样、交叉验证、预注册和外部复现都降低特定偏差。它们不回答“自然为何持续”的全局问题，却把不可操作的普遍担忧转为可检验的局部可靠性。局部改进不等于哲学问题已经消失，但能解释科学为何并非裸露地从几次观察跳到普遍规律。 “新谜题”提醒归纳依赖分类方式。有限数据同时符合许多延伸规则，我们选择绿色而非某个古怪谓词，不只是因为数据本身。语言实践、理论简洁性、因果结构和投射历史共同约束哪些性质可归纳。所谓相似性并非世界自动贴好的单一标签。

\minitopic{比较视角：思、习与自我校正}《孟子》中关于“心之官则思”的论述把反思能力与感官牵引区分，同时置于道德修养的整体语境\parencite{mencius2008}；《论语》则反复把学与思并置\parencite{analects2003}。这些材料可与当代自我知识、先验反思和认知训练对话，但不能简单译为“理性主义战胜经验主义”。 “思”不是脱离经验的纯演绎机器，“学”也不只是被动感官输入。二者涉及注意、复习、实践、与他人问答以及对行动的校正。比较视角提示，自知与理性能力可能通过习惯和关系培养，而不是一项始终透明的内在资源。 当代理论的贡献在于分层：当前体验权威、状态分类、原因解释、长期记忆和规范判断具有不同证据结构。古典修养论的贡献在于提醒，主体如何形成可靠注意和诚实自察本身是认识论问题。两方互相施压，而不必被合成为同一概念表。

\minitopic{综合案例：多年后的研究访谈}研究团队在事件十年后访问参与者。多人鲜明记得一次公开承诺，早期会议纪要却没有记录；其中三人的叙述用词高度相似，可能来自近年的纪念文章。团队不能只以“多数记得”或“文件没有”裁决。 编码层要问参与者当时是否在场和注意；保存层考察十年间反复讲述；提取层检查提问是否暗示；来源层区分亲历、听闻和后来阅读。会议纪要也要审查记录目的、停录区间和编辑权限。外部文档不是无条件真相机器。 研究者可建立时间线，优先收集最早日记和邮件，分开访谈，要求参与者报告信心与来源，并寻找只有亲历者可能知道的非公开细节。相似措辞若都可追溯到同一文章，不构成独立证言；细节差异也不必意味着核心事件为假，因为真实记忆会有正常遗忘。 最终结论可以分级：“承诺发生有中等支持；公开场合与确切措辞证据不足；多数回忆可能受后期资料影响。”这种结论同时尊重第一人称记忆与公共核验。认识上的暂停和分层置信不是失败，而是当不同来源无法完全决定时最准确的状态。

\begin{tcolorbox}[breakable,colback=white,colframe=blue!30!black,title={本章研读任务},fonttitle=\bfseries]
选择一项自己确信的旧记忆，分别记录编码、保存、提取与来源信息；找出一项可能的击败者和一项独立核验。再选一个先验候选命题，区分概念学习所需经验与最终确证，并说明其是否必然为真。
\end{tcolorbox}
